% ЗГЕНЕРОВАНО scripts/build_format_v2_chapter.py — не редагувати вручну.
% Джерело: editorial/cards/20.*.json, format-v2/rozdil-20.json, paraleli.json.

\naukchapter{20}{14}{Глава 20. Ніщо не постає}
\naukintro{Остання глава повторює форму девʼятнадцятої й доводить її до кінця: чотирнадцять разів питає де те й де те, і щоразу відповідає не переліком, а собою. Останній вірш книги обриває й самий перелік: та й навіщо багато слів — у мене не постає ніщо. Книга, що почалася питанням учня, закінчується тим, що питати більше нічого.}
\vsec{20.1}
\vdeva{क्व भूतानि क्व देहो वा । क्वेन्द्रियाणि क्व वा मनः । क्व शून्यं क्व च नैराश्यं । मत्स्वरूपे निरञ्जने ॥}
\vtranslit{ква бгутані ква дехо ва | квендріяні ква ва манах | ква шуням ква ча найрашям | матсварупе ніраньджане ||}
\vperek{Де першоелементи, де тіло, де чуття, де розум? Де порожнеча, де без-надійність — у моїй власній формі, що не знає плями?}
\vcom{Глава 20 складається майже цілком із питань kva, і це не пошук, а спосіб сказати, що переліченого немає де шукати. Перший вірш задає рамку: усі пари й категорії поставлено до однієї locative-форми — matsvarūpe, «у моїй власній формі». Зверніть увагу, що поряд із тілом і чуттями стоїть nairāśya, без-надійність: це не чеснота, а один із станів, який теж зникає. Слово ātman тут не вжите — є svarūpa, «власна форма»; ототожнювати їх можна лише як тлумачення.}
\vpar{Українська паралель (наша інтерпретація)}{Григорій Сковорода, «Наркіс. Розмова про те: Пізнай себе».}{Наша паралель: у Сковороди питання де тіло й де чуття теж обертається на питання про того, хто питає. Наше порівняння — типологічне.}
\vfig{ilustratsii/20-1-neti.pdf}{Першоелементи й розум відпали: лишається власна форма}{Схема наша, не історична ілюстрація.}

\vsec{20.2}
\vdeva{क्व शास्त्रं क्वात्मविज्ञानं । क्व वा निर्विषयं मनः । क्व तृप्तिः क्व वितृष्णत्वं । गतद्वन्द्वस्य मे सदा ॥}
\vtranslit{ква шастрам кватмавіджнянам | ква ва нірвішаям манах | ква тріптіх ква вітрішнатвам | ґатадвандвася ме сада ||}
\vperek{Де писання, де пізнання Я, де розум без об'єктів? Де вдоволеність, де свобода від спраги — у мене, в кого пари вже відпали назавжди?}
\vcom{Другий вірш ставить під питання й засоби пізнання: писання, пізнання Я, безоб'єктний розум. Це важливо для розуміння всієї глави — заперечення не обминає духовних досягнень, воно починає саме з них. tṛpti і vitṛṣṇatva названо поруч і досить несподівано: вдоволеність теж виявляється тим, чого немає де шукати. Термін toṣa, з яким tṛpti часто ототожнюють, тут не вжитий; це різні слова, і переклад їх не змішує.}
\vpar{Українська паралель (наша інтерпретація)}{Іван Вишенський, «Книжка» (кінець XVI ст.).}{Наша паралель: у Вишенського писання й пізнання теж втрачають вагу перед тим, у кого пари вже відпали. Наше порівняння — типологічне.}
\vfig{ilustratsii/20-2-neti.pdf}{Писання й пізнання Я відпали}{Схема наша, не історична ілюстрація.}

\vsec{20.3}
\vdeva{क्व विद्या क्व च वाविद्या । क्वाहं क्वेदं मम क्व वा । क्व बन्धः क्व च वा मोक्षः । स्वरूपस्य क्व रूपिता ॥}
\vtranslit{ква відя ква ча вавідя | квахам кведам мама ква ва | ква бандгах ква ча ва мокшах | сварупася ква рупіта ||}
\vperek{Де знання й де незнання, де «я», де «це», де «моє»? Де пута, де визволення — і де взагалі якась оформленість у того, хто сам є своєю формою?}
\vcom{Вірш перебирає основні опозиції будь-якого вчення — знання й незнання, я й це, моє, пута й визволення — і всі ставить під те саме питання. Порядок показовий: mokṣa з'являється не як відповідь, а як останній член пари з bandha. Остання пада замикає коло: якщо він є власна форма (svarūpa), то приписувати йому ще й оформленість (rūpitā) немає де. Саме тут свідки розходяться в одному слові — rūpitā / rūpatā, — і вибір між ними змінює відтінок, але не напрям думки.}
\vpar{Українська паралель (наша інтерпретація)}{Памфіл Юркевич, «Серце і його значення в духовному житті людини» (1860).}{Наша паралель: коли вірш питає, де я й де моє, він розводить того, хто говорить, і те, що він собі приписує. Юркевич розводить їх так само рішуче. Наше порівняння — за способом розрізняти.}
\vfig{ilustratsii/20-3-neti.pdf}{Я, моє, пута й визволення відпали}{Схема наша, не історична ілюстрація.}

\vsec{20.4}
\vdeva{क्व प्रारब्धानि कर्माणि । जीवन्मुक्तिर् अपि क्व वा । क्व तद् विदेहकैवल्यं । निर्विशेषस्य सर्वदा ॥}
\vtranslit{ква прарабдгані кармані | джіванмуктір апі ква ва | ква тад відехакайвалям | нірвішешася сарвада ||}
\vperek{Де вже запущені в дію діяння, де звільнення ще за життя? Де те безтілесне кайвальям — у того, хто завжди без ознак?}
\vcom{Тут заперечено не лише пута, а й обидві форми звільнення, відомі традиції: прижиттєве (jīvanmukti) і посмертне (videhakaivalya). Разом із ними зникає й prārabdha-karman — залишкова карма, якою зазвичай пояснюють, чому звільнений ще живе в тілі. Логіка вірша: усе це ознаки, а того, хто без ознак (nirviśeṣa), означити нічим. Зауважмо, що слово mukti в главах 16, 17 і 20 ужите лише двічі — тут і в 20.12, — і обидва рази в переліку того, чого немає.}
\vpar{Українська паралель (ПОТРЕБУЄ ПЕРЕВІРКИ)}{українська філософська проза, не визначено.}{ПОТРЕБУЄ ПЕРЕВІРКИ: питання про безтілесне звільнення найближче до української філософської прози, але конкретного твору ми не перевірили.}
\vfig{ilustratsii/20-4-neti.pdf}{Діяння й безтілесне звільнення відпали}{Схема наша, не історична ілюстрація.}

\vsec{20.5}
\vdeva{क्व कर्ता क्व च वा भोक्ता । निष्क्रियं स्फुरणं क्व वा । क्वापरोक्षं फलं वा क्व । निःस्वभावस्य मे सदा ॥}
\vtranslit{ква карта ква ча ва бгокта | нішкріям спхуранам ква ва | квапарокшам пхалам ва ква | ніхсвабгавася ме сада ||}
\vperek{Де діяч, де той, хто зазнає? Де бездіяльне сяйво? Де безпосереднє знання, де його плід — у мене, в кого немає власної природи?}
\vcom{Вірш забирає найстійкішу пару будь-якої системи — діяча й споживача — і додає до неї niṣkriya sphuraṇa, бездіяльне сяйво: те, що виявляється, не діючи. Це єдина в главі спроба назвати щось позитивне, і вона одразу поставлена під питання kva. Останнє слово — niḥsvabhāva, «без власної природи»; воно не означає порожнечі, а вказує, що жодна природа не може бути йому приписана. Термінів cit чи sākṣin тут немає, хоч sphuraṇa за змістом найближче до них.}
\vpar{Українська паралель (наша інтерпретація)}{Микола Зеров, «Камена» (1924).}{Наша паралель: у Зерова бездіяльне сяйво теж постає як повнота, а не як відсутність. Наше порівняння — типологічне.}
\vfig{ilustratsii/20-5-flame.pdf}{Бездіяльне сяйво без власної природи}{Схема наша, не історична ілюстрація.}

\vsec{20.6}
\vdeva{क्व लोकः क्व मुमुक्षुर् वा । क्व योगी ज्ञानवान् क्व वा । क्व बद्धः क्व च वा मुक्तः । स्वस्वरूपे ऽहम् अद्वये ॥}
\vtranslit{ква локах ква мумукшур ва | ква йоґі джнянаван ква ва | ква баддгах ква ча ва муктах | свасварупе ’хам адває ||}
\vperek{Де світ, де той, хто прагне звільнення, де йог чи знавець? Де зв'язаний, де звільнений — я у своїй власній формі, яка не знає двоїстості.}
\vcom{Уперше в главі з'являється ствердне речення: aham ... svasvarūpe advaye, «я — у своїй недвоїстій формі». Усе інше лишається питанням. Показово, що серед зниклих постатей є й muktaḥ, звільнений: у недвоїстості немає ні того, хто зв'язаний, ні того, хто звільнений. Слово advaya, «недвоїстий», — єдиний явний адвайтичний термін у цих трьох главах; усе решта сказано описово. Це місце варто читати уважно: «я» тут не протиставлене світові, воно вказане як те, у чому світу не знаходиться місця.}
\vpar{Українська паралель (наша інтерпретація)}{Тарас Шевченко, «Кобзар» (1840).}{Наша паралель: у Шевченка питання де звʼязаний, де звільнений теж звучить як докір мірі, а не як опис двох людей. Наше порівняння — типологічне, не про джерело.}
\vfig{ilustratsii/20-6-balance.pdf}{Дві шальки: світ і визволення — обидві порожні}{Схема наша, не історична ілюстрація.}

\vsec{20.7}
\vdeva{क्व सृष्टिः क्व च संहारः । क्व साध्यं क्व च साधनम् । क्व साधकः क्व सिद्धिर् वा । स्वस्वरूपे ऽहम् अद्वये ॥}
\vtranslit{ква сріштіх ква ча самхарах | ква садгям ква ча садганам | ква садгаках ква сіддгір ва | свасварупе ’хам адває ||}
\vperek{Де творення, де згортання, де мета, де засіб? Де той, хто практикує, де здобуття — я у своїй власній формі, яка не знає двоїстості.}
\vcom{Вірш знімає всю структуру практики: мету, засіб, виконавця й результат. Це найпряміше місце глави, де сказано, що шлях як такий не має опори. Друга пада повторює 20.6 дослівно — прийом навмисний: після кожного переліку читача повертають до того самого місця. Слово siddhi тут ужите в звичайному значенні «здобуття», а не як назва надприродних сил у йозі.}
\vpar{Українська паралель (наша інтерпретація)}{Богдан-Ігор Антонич, «Ротації» (1938).}{Наша паралель: в Антонича творення й згортання теж описані як рух самої природи, а не чиясь справа. Наше порівняння — за образом.}
\vfig{ilustratsii/20-7-motion.pdf}{Творення й згортання відпали}{Схема наша, не історична ілюстрація.}

\vsec{20.8}
\vdeva{क्व प्रमाता प्रमाणं वा । क्व प्रमेयं क्व च प्रमा । क्व किंचित् क्व न किंचिद् वा । सर्वदा विमलस्य मे ॥}
\vtranslit{ква прамата праманам ва | ква прамеям ква ча прама | ква кімчіт ква на кімчід ва | сарвада вімалася ме ||}
\vperek{Де той, хто пізнає, де засіб пізнання, де предмет, де саме пізнання? Де «щось», де «ніщо» — у мене, хто завжди без плями?}
\vcom{Тут розібрано повний апарат пізнання — чотири члени, які індійські школи сперечалися уточнювати. Текст не заперечує жодного з них окремо: він питає, де вони, коли немає місця, у якому вони могли б стояти. Пада kva kiṃcit kva na kiṃcid vā забирає й останню опору — саму пару «є / немає». Це важливо для розуміння всієї глави: заперечення не є твердженням порожнечі, воно знімає обидві можливості.}
\vpar{Українська паралель (ПОТРЕБУЄ ПЕРЕВІРКИ)}{українська релігійна філософія, не визначено.}{ПОТРЕБУЄ ПЕРЕВІРКИ: перелік пізнавача, засобу й предмета, після якого питати нічого, близький до української релігійної філософії, але конкретного місця ми не звірили.}
\vfig{ilustratsii/20-8-neti.pdf}{Пізнавач, предмет і пізнання відпали}{Схема наша, не історична ілюстрація.}

\vsec{20.9}
\vdeva{क्व विक्षेपः क्व चैकाग्र्यं । क्व निर्बोधः क्व मूढता । क्व हर्षः क्व विषादो वा । सर्वदा निष्क्रियस्य मे ॥}
\vtranslit{ква вікшепах ква чайкаґр'ям | ква нірбодгах ква мудгата | ква харшах ква вішадо ва | сарвада нішкріяся ме ||}
\vperek{Де розсіяність, де зосередженість, де несвідомість, де затьмареність? Де радість, де відчай — у мене, хто завжди бездіяльний?}
\vcom{Вірш повторює схему 18.10 — три пари, кожна заперечує обидва боки, — але з іншим закінченням: замість «у вгамованого йога» тут «у мене, хто завжди бездіяльний» (niṣkriya). niṣkriya не означає «лінивий»: це відсутність діяння як діяча, а не відсутність руху. Пари підібрані так, щоб охопити і стан уваги (розсіяність / зосередженість), і стан розуму (несвідомість / затьмареність), і стан почуття (радість / відчай). Те, що лишається поза цими парами, у вірші не названо.}
\vpar{Українська паралель (наша інтерпретація)}{Василь Стус, «Палімпсести».}{Наша паралель: у Стуса розсіяність і зосередженість теж не творять пари для того, хто вже стоїть. Наше порівняння — типологічне.}
\vfig{ilustratsii/20-9-balance.pdf}{Дві шальки: розсіяність і зосередженість}{Схема наша, не історична ілюстрація.}

\vsec{20.10}
\vdeva{क्व चैष व्यवहारो वा । क्व च सा परमार्थता । क्व सुखं क्व च वा दुःखं । निर्विमर्शस्य मे सदा ॥}
\vtranslit{ква чайша в'явахаро ва | ква ча са парамартхата | ква сукхам ква ча ва духкхам | нірвімаршася ме сада ||}
\vperek{Де ця буденна практика, де та вища істина? Де щастя, де страждання — у мене, хто вже завжди нічого не розглядає?}
\vcom{Найнебезпечніше місце для читача: тут знято пару vyavahāra / paramārtha, тобто буденний і вищий рівні істини. Зазвичай цю пару використовують, щоб пояснити, чому світ хибний, а Брахмо істинне; вірш не робить і цього кроку — обидва рівні однаково не мають де стояти. nirvimarśa, «без розгляду», — не заклик до нерозуміння, а опис: розглядати нічого, бо немає двох речей, які можна було б зіставити.}
\vpar{Українська паралель (наша інтерпретація)}{Ліна Костенко, «Берестечко» (1999).}{Наша паралель: у Костенко буденна практика й вища істина теж втрачають різницю перед обличчям того, що вже є. Наше порівняння — за будовою думки.}
\vfig{ilustratsii/20-10-balance.pdf}{Дві шальки: практика й вища істина}{Схема наша, не історична ілюстрація.}

\vsec{20.11}
\vdeva{क्व माया क्व च संसारः । क्व प्रीतिर् विरतिः क्व वा । क्व जीवः क्व च तद्ब्रह्म । सर्वदा विमलस्य मे ॥}
\vtranslit{ква мая ква ча самсарах | ква прітір віратіх ква ва | ква джівах ква ча тадбрахма | сарвада вімалася ме ||}
\vperek{Де майя, де сансара, де прихильність і де відвернення? Де жива істота, де те Брахмо — у мене, хто завжди без плями?}
\vcom{Вірш називає основні поняття будь-якого вчення про визволення — майю, сансару, прихильність і відвернення, живу істоту й Брахмо — і ставить їх в один ряд із питанням kva. Найрадикальніше тут те, що Брахмо стоїть у тому самому ряду, що й сансара: текст не протиставляє їх як хибне й істинне. Це місце легко прочитати як спростування адвайти; точніше сказати, що вірш відмовляється від її словника, не пропонуючи іншого.}
\vpar{Українська паралель (наша інтерпретація)}{Михайло Коцюбинський, «Тіні забутих предків» (1911).}{Наша паралель: у Коцюбинського мая й сансара теж постають як те, що показується й зникає, не питаючи людину. Наше порівняння — за образом.}
\vfig{ilustratsii/20-11-illusion.pdf}{Мая й сансара під серпанком}{Схема наша, не історична ілюстрація.}

\vsec{20.12}
\vdeva{क्व प्रवृत्तिर् निवृत्तिर् वा । क्व मुक्तिः क्व च बन्धनम् । कूटस्थनिर्विभागस्य । स्वस्थस्य मम सर्वदा ॥}
\vtranslit{ква правріттір нівріттір ва | ква муктіх ква ча бандганам | кутастханірвібгаґася | свастхася мама сарвада ||}
\vperek{Де діяння і де припинення, де звільнення, де пута — у мене, незмінного й неділимого, що завжди в собі?}
\vcom{Пару pravṛtti / nivṛtti вже було знято в 16.8, і тут вона повертається разом із парою mukti / bandhana. Друга половина вірша — єдина в главі, де подано ствердні означення: kūṭastha, nirvibhāga, svastha. Усі три описують не властивість, а відсутність поділу: незмінний, неділимий, у собі. Саме тому питання kva не має відповіді-місця — ділити нема на що.}
\vpar{Українська паралель (наша інтерпретація)}{Іван Франко, «Мойсей» (1905).}{Наша паралель: у Франка діяння й припинення теж не утворюють вибору для того, хто незмінний. Наше порівняння — типологічне.}
\vfig{ilustratsii/20-12-balance.pdf}{Дві шальки: діяння й припинення}{Схема наша, не історична ілюстрація.}

\vsec{20.13}
\vdeva{क्वोपदेशः क्व वा शास्त्रं । क्व शिष्यः क्व च वा गुरुः । क्व चास्ति पुरुषार्थो वा । निरुपाधेः शिवस्य मे ॥}
\vtranslit{квопадешах ква ва шастрам | ква шішях ква ча ва ґурух | ква часті пурушартхо ва | нірупадгех шівася ме ||}
\vperek{Де настановлення, де писання, де учень, де вчитель? Де й мета людського життя — у мене, хто без жодної опори, у мене, блаженного?}
\vcom{Після того як у 20.2 знято писання, тут знято всю установу передачі знання: настановлення, текст, учня, вчителя. Це послідовно: якщо немає ні знання, ні незнання, то немає й того, кому й від кого його передавати. puruṣārtha зникає останнім — разом із самою ідеєю мети. śiva тут майже напевно прикметник «блаженний, ласкавий», а не ім'я бога; вірш не дає приводу читати його як теологічне твердження, і переклад цього не робить.}
\vpar{Українська паралель (наша інтерпретація)}{Григорій Сковорода, «Алфавіт, чи Буквар миру» (діалог).}{Наша паралель: у Сковороди настановлення й писання теж відступають перед тим, хто вже не має опори. Наше порівняння — типологічне.}
\vfig{ilustratsii/20-13-neti.pdf}{Настановлення, учень і вчитель відпали}{Схема наша, не історична ілюстрація.}

\vsec{20.14}
\vdeva{क्व चास्ति क्व च वा नास्ति । क्वास्ति चैकं क्व च द्वयम् । बहुनात्र किम् उक्तेन । किंचिन् नोत्तिष्ठते मम ॥}
\vtranslit{ква часті ква ча ва насті | квасті чайкам ква ча дваям | бахунатра кім уктена | кімчін ноттіштхате мама ||}
\vperek{Де «є» і де «немає», де «одне» і де «двоє»? Та й навіщо багато слів: у мене не постає ніщо.}
\vcom{Останній вірш закриває главу відмовою від подальшої мови: після того як знято пару «є / немає» і пару «одне / двоє», говорити більше нема про що. Формула bahunātra kim uktena — усталений спосіб обірвати перелік, і вона тут доречна подвійно: як риторика і як зміст. Слово kiṃcit у заперечному реченні означає «ніщо»; це не твердження порожнечі, а констатація, що ніщо не підводиться як предмет. Цим і завершується глава: не відповіддю, а припиненням питання.}
\vpar{Українська паралель (наша інтерпретація)}{Олександр Довженко, «Зачарована Десна» (1956).}{Наша паралель: у Довженка останнє слово теж не додає нічого до сказаного, а лише закриває розмову. Наше порівняння — за будовою думки.}
\vfig{ilustratsii/20-14-space.pdf}{Ніщо не постає: останнє слово}{Схема наша, не історична ілюстрація.}

